If $Y$ is a line, take $U=(\mathbb{P}^2)^*-\{Y\}$. Otherwise, the closure of the image of the tangent morphism in \exref{I.7.3} has dimension at most $1$. The singular locus of $Y$ is finite, and the lines through each singular point form a line in $(\mathbb{P}^2)^*$. Their union with the closed tangent image is therefore a proper closed subset. Let $U$ be its complement. Every line in $U$ meets $Y$ only at nonsingular points and is transverse at each of them, so every local intersection multiplicity is $1$. By \exref{I.5.4}, their sum is $d$, and there are exactly $d$ intersection points.
